% interface=en

\environment mstytools

\setvariables[tools][banner=XML]

\starttext

\starttextmakeup

    \setupframedtexts[explanationtext][depthcorrection=no]

    \startexplanation
        \setups[banner]
    \stopexplanation

    \startexample
        This program, written in \RUBY, collects a few handy tools we use
        at \PRAGMA\ to generate and (pre)process \XML.

        Most commands operate on the given filename or, when no name is
        given, on all files that match a default pattern.

        \starttyping
xmltools --version
xmltools --recurse --dir --pattern=*.tex --output=textools.xml
        \stoptyping

        These examples show a few general options:

        \starttabulate[|Tlf{\blue\tt}|p|]
        \NC version \NC report banner         \NC \NR
        \NC recurse \NC recurse into subpaths \NC \NR
        \stoptabulate
    \stopexample

\stoptextmakeup

% directories

\starttextmakeup

    \startexplanation

        The \type {--dir} and \type {-ls} commands generate a file with
        filenames and file info, organized by path. The pattern and output
        file switches are mandate. You can specify an url, which will be
        added to the specification for later usage.

        \starttabulate[|Tlf{\blue\tt}|p|]
        \NC pattern   \NC the files to list           \NC \NR
        \NC output    \NC the output file             \NC \NR
        \NC url       \NC an url to add to the file   \NC \NR
        \NC stripname \NC strip basename from pattern \NC \NR
        \NC root      \NC the path to run on          \NC \NR
        \stoptabulate

        The files generated by this command can be processed by \CONTEXT\
        in several ways (see \type {x-dir-*.tex} files).

    \stopexplanation

    \startexample
        \starttyping
xmltools --dir --pattern=*.tex --output=textools.xml
xmltools --dir --recurse --root=e:/context/sources --pattern=*.tex --output=list.xml
        \stoptyping
    \stopexample

\stoptextmakeup

% analyze

\starttextmakeup

    \startexplanation

        The \type {--analyze} command results in a xml formatted log file
        describing what elements, attributes and entities are used.

        \starttabulate[|Tlf{\blue\tt}|p|]
        \NC pattern \NC the files to list \NC \NR
        \NC output  \NC the output file   \NC \NR
        \stoptabulate

        The files generated by this command can be processed by \CONTEXT\
        into a pretty print using (see \type {x-xml-11.tex} file).

        \starttyping
texexec --pdf --use=xml-analyze myfile.xlg
        \stoptyping

    \stopexplanation

    \startexample
        \starttyping
xmltools --analyze myfile.xml
xmltools --analyze --output=resultfile.xml myfile.xml
        \stoptyping
    \stopexample

\stoptextmakeup

\stoptext
